% sec-09: the eight-stage build and the five diagram vocabularies.  Figure 18.
\section{Build Method}

\subsection{Two Schedules, Thirteen Sub-Prompts}

The whole build runs from one master prompt, which generates thirteen stage
sub-prompts in two schedules that never share a stage: five drive the ten
application file sets, eight drive this paper. Each names its output directory,
its commit floor, and the repository files it may read.

\begin{apptable}
\begin{tabularx}{\textwidth}{>{\raggedright\arraybackslash}p{1.2cm} >{\raggedright\arraybackslash}p{0.7cm} >{\raggedright\arraybackslash}p{3.5cm} >{\raggedright\arraybackslash}X >{\raggedright\arraybackslash}p{1.5cm}}
\headrow \thead{Part} & \thead{No} & \thead{Stage} & \thead{Output} & \thead{Figures}\\
\midrule
PART I & 1 & Recipients and templates & Recipient list, cover variants, style and bib contract & none \\
PART I & 2 & Set A, surgical & Applications 01 to 05 & 17 \\
PART I & 3 & Set B, medical oncology & Applications 06 to 10 & 17 \\
PART I & 4 & Email and attachments & Ten `.txt` files, ten Overleaf zips & none \\
PART I & 5 & READMEs and audit & Every `funding/**` README; the error-fix pass & none \\
PART II & 1 & mermaid & Six figure specifications & 6 \\
PART II & 2 & plantuml & Three figure specifications & 3 \\
PART II & 3 & d2 & Four figure specifications & 4 \\
PART II & 4 & diagrams (python) & Three figure specifications & 3 \\
PART II & 5 & graphviz & Four figure specifications & 4 \\
PART II & 6 & draft-apply & Skeleton, drafting instructions, 20 figure slots & 20 placed \\
PART II & 7 & full-apply & Every instruction resolved, every figure drawn & 20 drawn \\
PART II & 8 & final-apply & Senior-author proof-reading pass & 20 polished \\
\end{tabularx}
\end{apptable}
\tabcap{The thirteen stage sub-prompts in two schedules. The application stages
carry 34 figures across ten attachments; the paper stages carry 20, specified
once and drawn twice.}

\subsection{Why the Diagram Split Is Uneven}

An equal split across five platforms would be four figures each. The split used
here is 6, 3, 4, 3, 4, and the difference is not arbitrary.

\begin{apptable}
\begin{tabularx}{\textwidth}{>{\raggedright\arraybackslash}p{3.5cm} >{\raggedright\arraybackslash}p{1.2cm} >{\raggedright\arraybackslash}X}
\headrow \thead{Platform} & \thead{Count} & \thead{The question only this vocabulary can state}\\
\midrule
Mermaid-type & 6 & What happens next, what decides, how long, and who spoke in what order \\
PlantUML-type & 3 & Who is permitted to act, under what guard a state changes, and what runs concurrently \\
D2-type & 4 & What contains what, what tabulates against what, and what joins to what \\
Diagrams (python)-type & 3 & What runs where, across which trust boundary, and on whose hardware \\
Graphviz-type & 4 & What depends on what, how a failure propagates, and what a single record answers \\
\end{tabularx}
\end{apptable}
\tabcap{The five platforms, their share of the twenty figures, and the question
each is the only one able to state. Forcing an equal split would have moved at
least four figures into a vocabulary that cannot carry their subject.}

Mermaid takes the largest share because the paper's spine is chronological: a
policy document, ten submissions, a fourteen-month record, and one operative day
are all sequences. PlantUML takes the smallest because only three subjects need
a guard grammar. Forced to four each, the authority diagram would have become a
flowchart, which cannot draw the absence of a permission; the evidence tiers a
ladder, which asserts order where the paper means containment; and the fault
tree a dependency graph, which cannot distinguish AND from OR.

\subsection{One Operator Against a Program Office}

\begin{appfloat}[!tb]
\begin{appfig}[diagrams (python)-type / clustered by layer]
% Two clusters at x = 0 and x = 7.4, five tiles each at the same y positions so
% layers align horizontally and a row is a like-for-like comparison.  Cost notes
% sit to the right of the right-hand column, outside both clusters.
\foreach \i/\y/\g/\lab in {1/0/{\glyphdoc}/{Literature: 40 meta-analyses},
  2/-2.1/{\glyphchart}/{Modeling: 3 simulations}, 3/-4.2/{\glyphshield}/{Regulatory: protocol and IND},
  4/-6.3/{\glyphcloud}/{Release: DOI and repository}, 5/-8.4/{\glyphflask}/{Clinical: contracted site}}{
  \dgnode{dgtile}{0}{\y}{\g}{\lab}{p\i}}
\begin{scope}[on background layer]
\node[dgcluster,fit=(p1)(p1l)(p5)(p5l),inner sep=7pt] (cl) {};
\end{scope}
\node[dgctitle,anchor=south west] at ([yshift=1mm]cl.north west) {One independent scientist, 14 months};
\foreach \i/\y/\g/\lab in {1/0/{\glyphteam}/{Scientific writing group},
  2/-2.1/{\glyphcpu}/{Modeling and simulation team}, 3/-4.2/{\glyphgear}/{Regulatory affairs unit},
  4/-6.3/{\glyphdb}/{Data management office}, 5/-8.4/{\glyphuser}/{In-house trial unit}}{
  \dgnodeg{dgtileg}{7.4}{\y}{\g}{\lab}{q\i}}
\begin{scope}[on background layer]
\node[dgcluster2,fit=(q1)(q1l)(q5)(q5l),inner sep=7pt] (cr) {};
\end{scope}
\node[dgctitle2,anchor=south west] at ([yshift=1mm]cr.north west) {The equivalent institutional program office};
\foreach \i in {1,...,4}{\pgfmathtruncatemacro{\j}{\i+1}
  \draw[dgedgeb] (p\i l.south) -- (p\j.north);
  \draw[dgedged] (q\i l.south) -- (q\j.north);}
\node[dgtilem,minimum width=9mm,minimum height=9mm] (sh) at (3.7,-8.4) {};
\glyphlink{sh}{protowhite}
\node[dglabel,anchor=north] at ([yshift=-5.4mm]sh.center) {shared layer};
\draw[dgbi] (p5) -- (sh);
\draw[dgbi] (sh) -- (q5);
\pnote{9.6}{-2.1}{\$36{,}330 per run\\against \$120{,}000 to \$2M}
\pnote{9.6}{-4.2}{no indirect cost against\\50 to 60 percent}
\node[font=\tiny\itshape,text=protogray,anchor=north,text width=104mm,align=center]
  at (3.7,-9.9) {Layers align horizontally, so a row is a like-for-like function
  comparison and a column is one organization's whole stack. Only the bottom
  layer is shared, and it is the one neither can do alone.};
\end{appfig}
\vspace{-0.7cm}
\figcaption{Five layers, executed twice. The left column is one operator with an
LLM toolchain and a public repository; the right is the institutional
equivalent. Only the bottom layer is common to both.}
\end{appfloat}

\subsection{What Is Reproducible, and What Is Not}

Reproducible: the master prompt, the thirteen sub-prompts, the named sources,
the commit sequence, the figure construction notes, and the compile command. A
reader with the repository can rebuild every artifact and check every number
against the file it came from.

Not reproducible: the judgement at each step. Which recipient qualifies, which
vocabulary a figure needs, which sentence in a 123-page report is worth quoting,
and which of the author's own claims is wrong, these are decisions, and a second
run would make some differently. That is why the author's review is a required
stage rather than an optional one, and why the PART I audit pass, which found
the first-in-human error in nine of ten applications, is recorded in \S7 rather
than absorbed silently.
